% from ega3-1-fr.tex
\oldpage[0\textsubscript{III}]{10}
objet de $C$ muni d'une loi de composition interne $\gamma_{X''}$, et
$w':X'\to X''$ un morphisme dans $C$ qui est un homomorphisme pour
$\gamma_{X'}$ et $\gamma_{X''}$, il est clair que le morphisme
$w'w:X\to X''$ est un homomorphisme pour les lois de composition $\gamma_X$
et $\gamma_{X''}$. Un isomorphisme $w:X\xrightarrow{\sim}X'$ dans $C$ est
appelé \emph{isomorphisme pour les lois de composition $\gamma_X$ et
$\gamma_{X'}$} si $w$ est un homomorphisme pour ces lois de composition, et
si son morphisme réciproque $w^{-1}$ est un homomorphisme pour les lois de
composition $\gamma_{X'}$ et $\gamma_X$.

On définit de la même manière les \emph{dihomomorphismes} pour les couples
d'objets de $C$ munis de lois de composition externes.
\end{env}

\begin{env}[8.2.3]
\label{0.8.2.3-fr}
Lorsqu'une loi de composition interne $\gamma_X$ sur un objet $X\in C$ est
telle que $\gamma_X(Y)$ soit une \emph{loi de groupe} sur $h_X(Y)$ pour
\emph{tout} $Y\in C$, on dit que $X$, muni de cette loi, est un
\emph{$C$-groupe} ou un \emph{$C$-objet en groupes}. On définit de même les
\emph{$C$-anneaux}, \emph{$C$-modules}, etc.
\end{env}

\begin{env}[8.2.4]
\label{0.8.2.4-fr}
Supposons que le \emph{produit} $X\times X$ d'un objet $X\in C$ par lui-même
existe dans $C$~; par définition, on a alors
$h_{X\times X}=h_X\times h_X$ à un isomorphisme canonique près, puisqu'il
s'agit d'un cas particulier de limite projective
(\hyperref[0.8.1.9-fr]{8.1.9})~; une loi de composition interne sur $X$ peut
donc être considérée comme un morphisme fonctoriel
$\gamma_X:h_{X\times X}\to h_X$, et détermine donc canoniquement
(\hyperref[0.8.1.6-fr]{8.1.6}) un élément
$c_X\in\operatorname{Hom}(X\times X,X)$ tel que $h_{c_X}=\gamma_X$~; dans
ce cas, la donnée d'une loi de composition interne sur $X$ est donc
équivalente à celle d'un morphisme $X\times X\to X$~; lorsque $C$ est la
catégorie $\mathrm{Ens}$, on retrouve la notion classique de loi de
composition interne sur un ensemble. On a un résultat analogue pour une loi
de composition externe lorsque le produit $Z\times X$ existe dans $C$.
\end{env}

\begin{env}[8.2.5]
\label{0.8.2.5-fr}
Avec les notations précédentes, supposons en outre que $X\times X\times X$
existe dans $C$~; la caractérisation du produit comme objet représentant un
foncteur (\hyperref[0.8.1.9-fr]{8.1.9}) entraîne l'existence d'isomorphismes
canoniques
\[
  (X\times X)\times X\xrightarrow{\sim}X\times X\times X
  \xrightarrow{\sim}X\times(X\times X)~;
\]
si on identifie canoniquement $X\times X\times X$ à $(X\times X)\times X$,
l'application $\gamma_X(Y)\times 1_{h_X(Y)}$ s'identifie à
$h_{c_X\times 1_X}(Y)$ pour tout $Y\in C$. Il est par suite équivalent de
dire que pour tout $Y\in C$, la loi interne $\gamma_X(Y)$ est
\emph{associative}, ou que le diagramme d'applications
\[
  \xymatrix{
    h_X(Y)\times h_X(Y)\times h_X(Y)
      \ar[rr]^{\gamma_X(Y)\times 1}\ar[dd]_{1\times\gamma_X(Y)} &&
    h_X(Y)\times h_X(Y)\ar[dd]^{\gamma_X(Y)}\\\\
    h_X(Y)\times h_X(Y)\ar[rr]_{\gamma_X(Y)} && h_X(Y)
  }
\]
